\chapter*{About this book}
\markboth{About this book}{About this book}
\phantomsection
\addcontentsline{toc}{chapter}{About this book}
\label{um:ch:about}

\section*{Purpose of this book}
This book is the User Manual. It shows how to use the project: how to set up the tools, run the flows, write the code and write the books. It holds tasks and commands. It does not hold theory of operation. The theory is in the Technical Manual.
\input{about-books}

\section*{Intended audience}
The book is for every person who works in the repository: designers, verification engineers, physical design engineers, writers and reviewers.

\section*{Using this book}
Each chapter is a guide or a set of tasks. Use the table of contents to find one. The chapter on the documentation system explains the books, the label scheme and the writing rules. A task is a numbered list of steps. A command is in monospace text.

\input{about-conventions}

\section*{Feedback}
Open an issue in the project tracker. Name the label of the section, for example \code{um:ch:about}.
